\chapter{Maximum Likelihood (MLE) and Maximum a Posteriori (MAP)}
\companion{StatQuest: Maximum Likelihood, Clearly Explained}{\ytid{XepXtl9YKwc}}

MLE and MAP are named in the syllabus and quietly sit under almost everything else: least squares, log loss, cross-entropy,
Naive Bayes counts, Laplace smoothing, Ridge and Lasso, EM for mixtures. If you can explain where a loss function
\emph{comes from}, interviewers see that you understand ML rather than having memorised it.

\section{Probability vs likelihood: flipping the question}
\begin{itemize}
\item \textbf{Probability}: the parameter is known; how likely is the data? ``If the coin is fair ($p = 0.5$), what is the chance of
  7 heads in 10 tosses?'' Answer: $\binom{10}{7}0.5^{10} = 0.117$.
\item \textbf{Likelihood}: the data are known; how plausible is each parameter value? ``I saw 7 heads in 10 tosses. How well does
  $p = 0.5$ explain that? $p = 0.7$? $p = 0.9$?''
\end{itemize}
The formula is the same, $L(p) = P(\text{data}\mid p)$, but we treat it as a function of $p$ with the data fixed.

\begin{example}[: compare candidate parameters]
7 heads in 10. $L(p)\propto p^7(1-p)^3$. $L(0.5) = 0.00098$; $L(0.7) = 0.00222$; $L(0.9) = 0.00048$. The data are best explained
by $p = 0.7$. Note that the likelihood values do not sum to 1 over $p$: likelihood is not a probability distribution over parameters.
\end{example}

\section{Maximum likelihood estimation: the recipe}
\textbf{MLE} chooses the parameter that makes the observed data most probable: $\hat\theta = \argmax_\theta L(\theta)$.
\begin{enumerate}
\item Write the probability of one observation given $\theta$.
\item Assume observations are independent: the likelihood is the product over observations.
\item Take the log (products become sums; the maximiser is unchanged because log is increasing; numerically stable).
\item Differentiate with respect to $\theta$, set to zero, solve. (Check it is a maximum.)
\end{enumerate}

\subsection{Coin (Bernoulli)}
$k$ heads in $n$ tosses: $\ell(p) = k\log p + (n-k)\log(1-p)$. $\ell'(p) = \frac kp - \frac{n-k}{1-p} = 0 \Rightarrow \hat p = \frac kn$.
The MLE is the observed fraction. 12 heads in 20 tosses: $\hat p = 0.6$.

\subsection{Gaussian}
Data $x_1,\dots,x_n$ from $\mathcal N(\mu,\sigma^2)$:
$\ell = -\frac n2\log(2\pi\sigma^2) - \frac{1}{2\sigma^2}\sum(x_i - \mu)^2$.
\begin{itemize}
\item $\partial\ell/\partial\mu = 0 \Rightarrow \hat\mu = \bar x$ (the sample mean).
\item $\partial\ell/\partial\sigma^2 = 0 \Rightarrow \hat\sigma^2 = \frac1n\sum(x_i - \bar x)^2$: divides by $n$, not $n-1$, so it is
  \emph{biased} (slightly too small); the unbiased sample variance divides by $n-1$.
\end{itemize}
\begin{example}[: Gaussian MLE]
Data 2, 4, 6, 8: $\hat\mu = 5$, $\hat\sigma^2 = (9 + 1 + 1 + 9)/4 = 5$ (unbiased version: $20/3 = 6.67$).
\end{example}

\subsection{Poisson and exponential}
Poisson counts: $\hat\lambda = \bar x$. Applications per hour 3, 5, 4, 0, 3: $\hat\lambda = 3$.
Exponential waiting times with rate $\lambda$: $\hat\lambda = 1/\bar x$. Times 2, 3, 5 minutes: $\bar x = 3.33$, $\hat\lambda = 0.3$/min.

\subsection{A case where calculus fails: Uniform$(0,\theta)$}
$L(\theta) = \theta^{-n}$ if $\theta\ge\max x_i$, else 0. It decreases in $\theta$, so the MLE is the smallest allowed value:
$\hat\theta = \max x_i$. Data 1.2, 3.4, 2.7: $\hat\theta = 3.4$ (biased low: the true $\theta$ is almost surely a bit bigger).

\section{Why ML losses are negative log-likelihoods}
This is the single most useful idea in the chapter.
\begin{itemize}
\item \textbf{Regression with Gaussian noise}: $y = f(\mathbf x) + \varepsilon$, $\varepsilon\sim\mathcal N(0,\sigma^2)$.
  $-\log L = \frac{1}{2\sigma^2}\sum(y_i - f(\mathbf x_i))^2 + \text{const}$. \textbf{Minimising squared error = MLE.}
\item \textbf{Regression with Laplace noise}: $-\log L\propto\sum|y_i - f(\mathbf x_i)|$. \textbf{MAE = MLE} under Laplace noise.
\item \textbf{Binary classification}: $y\sim\text{Bernoulli}(p)$. $-\log L = -\sum[y\log p + (1-y)\log(1-p)]$.
  \textbf{Log loss = MLE.}
\item \textbf{Multiclass}: categorical distribution. \textbf{Cross-entropy = MLE.}
\item \textbf{Counts}: Poisson. \textbf{Poisson deviance = MLE} (Chapter 4).
\end{itemize}
So choosing a loss is choosing a noise model. Minimising cross-entropy between the data distribution and the model is also the same as
minimising the \textbf{KL divergence} from data to model, since the data's own entropy is a constant.

\section{MAP: adding prior beliefs}
\subsection{Bayes' theorem for parameters}
\[ \underbrace{P(\theta\mid D)}_{\text{posterior}} = \frac{\overbrace{P(D\mid\theta)}^{\text{likelihood}}\ \overbrace{P(\theta)}^{\text{prior}}}{P(D)}. \]
The \textbf{prior} encodes what we believe before seeing data. The \textbf{maximum a posteriori} estimate is the posterior's peak:
\[ \hat\theta_{MAP} = \argmax_\theta\big[\log P(D\mid\theta) + \log P(\theta)\big]. \]
It is MLE plus a term that pulls toward plausible values: a regulariser.

\subsection{Coin with a Beta prior}
A Beta$(a, b)$ prior on $p$ acts like $a - 1$ imaginary heads and $b - 1$ imaginary tails. With $k$ heads in $n$:
\[ \hat p_{MAP} = \frac{k + a - 1}{n + a + b - 2},\qquad \text{posterior mean} = \frac{k + a}{n + a + b}. \]
\begin{example}[: three heads in three tosses]
MLE: $3/3 = 1$: ``this coin never lands tails''. Absurd after three tosses. With a Beta(2,2) prior (mild belief in fairness):
MAP $= (3+1)/(3+2) = 0.8$; posterior mean $= 5/7 = 0.714$. With more data the prior's influence fades.
\end{example}

\begin{example}[: click-through rate smoothing]
A new job got 1 click in 2 views: MLE CTR $= 50\%$. Ranking by raw CTR would put it on top of jobs with 300 clicks in 10{,}000 views
(3\%). Use a Beta prior matching the platform average 5\%, e.g. Beta(2, 38): posterior mean $= (1 + 2)/(2 + 40) = 7.1\%$. Small
samples are pulled toward the average; large samples keep their own rate.
\end{example}

\subsection{Regularisation is MAP}
\begin{itemize}
\item Gaussian prior $w_j\sim\mathcal N(0,\tau^2)$: $-\log P(\mathbf w) = \frac{1}{2\tau^2}\norm{\mathbf w}^2 + c$. With Gaussian noise
  $\sigma^2$, MAP minimises $\frac{1}{2\sigma^2}\text{SSE} + \frac{1}{2\tau^2}\norm{\mathbf w}^2$, i.e. \textbf{Ridge} with
  $\lambda = \sigma^2/\tau^2$.
\item Laplace prior $p(w_j)\propto e^{-|w_j|/b}$: \textbf{Lasso}. The Laplace density has a sharp peak at 0, which is why it produces exact
  zeros.
\item Laplace smoothing in Naive Bayes: MAP (or posterior mean) with a Dirichlet prior.
\end{itemize}
\begin{example}[: Ridge's $\lambda$ from priors]
Noise variance $\sigma^2 = 4$, prior variance $\tau^2 = 2$: $\lambda = 4/2 = 2$.
\end{example}

\subsection{Conjugate priors}
When prior and posterior are in the same family, updating is just adding counts:
Beta--Bernoulli/Binomial, Gamma--Poisson, Dirichlet--Multinomial, Gaussian--Gaussian (known variance).
\begin{example}[: Gamma--Poisson]
Prior Gamma$(\alpha = 2, \beta = 1)$ on a rate; observe counts 3 and 5. Posterior Gamma$(2 + 8, 1 + 2) = $ Gamma$(10, 3)$; mean $10/3 = 3.33$.
\end{example}
\begin{example}[: Gaussian--Gaussian]
Prior $\mu\sim\mathcal N(0, 1)$, one observation $x = 4$ with noise variance 1: posterior mean $= \frac{0/1 + 4/1}{1/1 + 1/1} = 2$, a
precision-weighted average of prior mean and data.
\end{example}

\section{MLE vs MAP vs full Bayesian inference}
\begin{itemize}
\item \textbf{MLE}: one best value, no prior. Can overfit with little data (the 3/3 coin).
\item \textbf{MAP}: one best value, with prior. Equivalent to regularised MLE. Not invariant to reparameterisation (the peak of a
  density moves when you change variables).
\item \textbf{Full Bayes}: keep the whole posterior distribution and average predictions over it: $P(y\mid x, D) = \int P(y\mid x,\theta)
  P(\theta\mid D)\,d\theta$. Gives uncertainty; usually expensive (MCMC, variational inference).
\end{itemize}
As $n\to\infty$, all three agree (the likelihood dominates the prior).

\section{Properties of MLE worth quoting}
Under regularity conditions MLE is \textbf{consistent} (converges to the truth), \textbf{asymptotically normal}, and
\textbf{asymptotically efficient} (achieves the Cram\'er--Rao lower bound: variance $\approx 1/(nI(\theta))$ with Fisher information
$I$). It is \textbf{invariant}: the MLE of $g(\theta)$ is $g(\hat\theta)$ (MLE of $\sigma$ is $\sqrt{\hat\sigma^2}$). It can be biased in small
samples ($\hat\sigma^2$, Uniform). Bernoulli Fisher information is $\frac{1}{p(1-p)}$, giving standard error $\sqrt{p(1-p)/n}$.

\section{When there is no closed form: EM}
With hidden (latent) variables (which mixture component generated each point), the likelihood is a sum inside a log and has no closed
form. The \textbf{Expectation--Maximisation} algorithm alternates: E-step: compute the probability of each hidden assignment given
current parameters; M-step: re-estimate parameters as if those soft assignments were observed. Each iteration never decreases the
likelihood. Chapter 15 works through EM for Gaussian mixtures.

\begin{practical}[: MLE and MAP at InfoEdge]
\begin{itemize}
\item \textbf{CTR/apply-rate smoothing} for new jobs and new recruiters (Beta priors) before ranking.
\item \textbf{Salary engine}: shrink company-specific salary estimates toward the industry average when a company has few data points
  (hierarchical/empirical Bayes = MAP with a learned prior).
\item ``Bayesian inference'' is on the Tools slide: A/B test analysis with Beta posteriors (``probability variant B is better'').
\end{itemize}
\end{practical}

\stuck{StatQuest: Probability is not Likelihood}{\ytid{pYxNSUDSFH4}}
\section{Interview questions}

\begin{qa}{What is the difference between probability and likelihood?}
\oneline{Probability fixes the parameters and asks how likely different data are; likelihood fixes the observed data and asks how well different
parameter values explain it; same formula, different variable.}
\why Probability sums to 1 over outcomes; likelihood does not sum or integrate to 1 over parameters. Worked example: $L(0.7) > L(0.5) > L(0.9)$ for
7/10 heads.
\end{qa}

\begin{qa}{\pyq{OA style: cross-entropy as likelihood} Show that minimising log loss is maximum likelihood, and that minimising squared error is
maximum likelihood under Gaussian noise.}
\oneline{For Bernoulli labels the negative log-likelihood is exactly the log loss; for $y = f(x) + $ Gaussian noise, the negative log-likelihood is
$\frac{1}{2\sigma^2}\sum(y - f(x))^2$ plus a constant.}
\why Bernoulli: $P(y\mid p) = p^y(1-p)^{1-y}$; $-\log = -[y\log p + (1-y)\log(1-p)]$; summed over independent examples. Gaussian:
$P(y\mid x) = \frac{1}{\sqrt{2\pi}\sigma}e^{-(y - f)^2/(2\sigma^2)}$; $-\log = \frac{(y-f)^2}{2\sigma^2} + \log(\sqrt{2\pi}\sigma)$. Since $\sigma$ does not
depend on $f$, minimising SSE maximises likelihood.
\pushes \emph{``So what does choosing MAE assume?''} Laplace noise; its MLE location is the median.
\end{qa}

\begin{qa}{\pyq{INT: Bayes in real life; Naive Bayes smoothing as MAP} Explain Laplace smoothing as a Bayesian (MAP) estimate.}
\oneline{Adding 1 to every count is the posterior estimate under a uniform Dirichlet prior over word probabilities: it is as if we had seen every word once
in every class before looking at the data.}
\why Dirichlet$(\alpha_1,\dots,\alpha_V)$ prior with multinomial counts $c_w$ gives a Dirichlet$(c_w + \alpha_w)$ posterior; its mean is
$\frac{c_w + \alpha}{N + V\alpha}$. With $\alpha = 1$ that is Laplace smoothing (posterior mean; the MAP with $\alpha = 2$ gives the same formula with
$\alpha - 1 = 1$). It prevents zero probabilities and shrinks rare-word estimates toward uniform.
\end{qa}

\begin{qa}{Derive the MLE of $p$ for $k$ successes in $n$ Bernoulli trials.}
\oneline{$\hat p = k/n$.}
\why $\ell(p) = k\log p + (n-k)\log(1-p)$; $\ell' = k/p - (n-k)/(1-p) = 0\Rightarrow k(1-p) = (n-k)p\Rightarrow p = k/n$; $\ell'' < 0$.
\end{qa}

\begin{qa}{Derive the MLE of $\mu$ and $\sigma^2$ for a Gaussian. Is $\hat\sigma^2$ unbiased?}
\oneline{$\hat\mu = \bar x$ and $\hat\sigma^2 = \frac1n\sum(x_i - \bar x)^2$, which is biased low by the factor $(n-1)/n$.}
\why The bias arises because $\bar x$ is fitted to the same data, so deviations from $\bar x$ are smaller on average than deviations from the true $\mu$.
Dividing by $n-1$ (Bessel's correction) fixes it. Worked: data 2, 4, 6, 8 give 5 and 5 (unbiased 6.67).
\end{qa}

\begin{qa}{Three heads in three tosses. Give the MLE, the MAP with a Beta(2,2) prior and the posterior mean. Which would you use?}
\oneline{MLE 1.0, MAP 0.8, posterior mean 0.714; with three tosses I would use the posterior mean or MAP because the MLE overfits to tiny data.}
\end{qa}

\begin{qa}{\deep Show that L2 regularisation is MAP with a Gaussian prior and derive $\lambda$.}
\oneline{With Gaussian noise $\sigma^2$ and prior $w_j\sim\mathcal N(0,\tau^2)$, the negative log-posterior is $\frac{1}{2\sigma^2}\norm{\mathbf y - X\mathbf w}^2 +
\frac{1}{2\tau^2}\norm{\mathbf w}^2$; multiplying by $2\sigma^2$ gives Ridge with $\lambda = \sigma^2/\tau^2$.}
\why Intuition: a strong prior (small $\tau^2$) or noisy data (large $\sigma^2$) means large $\lambda$: trust the prior more. $\sigma^2 = 4$, $\tau^2 = 2$
gives $\lambda = 2$.
\pushes \emph{``And L1?''} Laplace prior. \emph{``Is the Lasso solution the posterior mean?''} No, it is the posterior mode; the posterior mean under a
Laplace prior is not sparse.
\end{qa}

\begin{qa}{\deep A new job has 1 click from 2 impressions. How would you estimate its CTR for ranking, and why not just use 50\%?}
\oneline{Use a Beta prior fitted to the platform's CTR distribution and take the posterior mean $(k + \alpha)/(n + \alpha + \beta)$; raw 50\% from 2 impressions is
extremely noisy and would wrongly rank the job above proven performers.}
\why With Beta(2, 38) (mean 5\%): $(1+2)/(2+40) = 7.1\%$. Fit $\alpha,\beta$ by method of moments on historic jobs (empirical Bayes). For exploration,
sample from the posterior (Thompson sampling) instead of using the mean.
\end{qa}

\begin{qa}{What is the MLE of $\theta$ for Uniform$(0,\theta)$ data 2.1, 4.7, 3.9, 1.2? Why can't you use calculus?}
\oneline{$\hat\theta = 4.7$; the likelihood $\theta^{-n}$ is decreasing on the allowed region $\theta\ge\max x_i$ and zero below it, so the maximum is at the
boundary, where the derivative is not zero.}
\end{qa}

\begin{qa}{What is the MLE of $\lambda$ for Poisson counts 3, 5, 4, 0, 3, and for an exponential with waiting times 2, 3, 5?}
\oneline{Poisson: $\bar x = 3$; exponential: $1/\bar x = 0.3$ per minute.}
\why Poisson: $\ell = \sum[x_i\log\lambda - \lambda] + c$; $\ell' = \sum x_i/\lambda - n = 0$. Exponential: $\ell = n\log\lambda - \lambda\sum x_i$;
$\ell' = n/\lambda - \sum x_i = 0$.
\end{qa}

\begin{qa}{Why do we maximise the log-likelihood instead of the likelihood?}
\oneline{Log turns the product over observations into a sum (easier to differentiate and to optimise with SGD), avoids numerical underflow, and does not move
the maximiser because log is strictly increasing.}
\end{qa}

\begin{qa}{Explain MLE vs MAP vs full Bayesian inference in three lines.}
\oneline{MLE: the parameter that best explains the data; MAP: the parameter that best explains the data and the prior (regularised MLE); full Bayes: keep the
whole posterior and average predictions over it, which also gives uncertainty.}
\end{qa}

\begin{qa}{\deep Why is MAP not invariant to reparameterisation while MLE is?}
\oneline{A likelihood is a function of the parameter, so its maximiser maps through any reparameterisation; a posterior is a density, and changing variables
multiplies it by a Jacobian, which can move its peak.}
\why If $\phi = g(\theta)$, $p_\phi(\phi) = p_\theta(g^{-1}(\phi))|dg^{-1}/d\phi|$. The extra factor shifts the mode. Example: the mode of a posterior over $\sigma$
differs from the square root of the mode over $\sigma^2$.
\end{qa}

\begin{qa}{What does minimising cross-entropy have to do with KL divergence?}
\oneline{Cross-entropy $H(p, q) = H(p) + KL(p\,\|\,q)$; since the data entropy $H(p)$ does not depend on the model, minimising cross-entropy minimises the KL
divergence from the data distribution to the model.}
\end{qa}

\begin{qa}{What are conjugate priors? Give three pairs and a worked update.}
\oneline{A prior is conjugate to a likelihood if the posterior stays in the prior's family: Beta--Binomial, Gamma--Poisson, Dirichlet--Multinomial; e.g.
Gamma(2,1) prior and counts 3, 5 give Gamma(10, 3), mean 3.33.}
\end{qa}

\begin{qa}{\deep Logistic regression on perfectly separable data has no MLE. How does a prior fix this?}
\oneline{The likelihood keeps increasing as $\norm{\mathbf w}\to\infty$, so no maximiser exists; a Gaussian prior adds $-\norm{\mathbf w}^2/(2\tau^2)$ to the log-posterior,
which eventually outweighs the shrinking gains in likelihood, so a finite MAP exists.}
\end{qa}

\begin{qa}{What is Fisher information and what does it tell you?}
\oneline{It measures how sharply the log-likelihood is curved around the true parameter, i.e. how much information each observation carries; the MLE's variance
is approximately $1/(nI(\theta))$.}
\worked Bernoulli: $I(p) = 1/(p(1-p))$; with $p = 0.1$, $n = 900$: SE $\approx\sqrt{0.09/900} = 0.01$.
\end{qa}

\begin{qa}{\deep In a salary engine, a small startup has 3 salary records for ``Data Scientist''. How do you estimate its typical salary?}
\oneline{Use a hierarchical (empirical Bayes) estimate that shrinks the company's mean toward the industry/city mean in proportion to how few records it has:
$\hat\mu_c = \frac{n_c\bar x_c/\sigma^2 + \mu_0/\tau^2}{n_c/\sigma^2 + 1/\tau^2}$.}
\why This is the Gaussian--Gaussian posterior mean. With $n_c = 3$ the industry mean carries real weight; with $n_c = 300$ the company's own mean dominates.
It is also how mixed-effects models and target encoding with priors work.
\end{qa}

\begin{qa}{How is MLE used to train neural networks?}
\oneline{Training with cross-entropy or MSE is maximising the likelihood of the labels under the network's output distribution; weight decay adds a Gaussian prior
(MAP); and the output layer (sigmoid, softmax, linear) defines which distribution is assumed.}
\end{qa}
